\begin{enumerate}[label=\textbf{1.1.\arabic*.},ref=1.1.\arabic*]
\setlength{\itemsep}{0.35em}
\item \label{formula-1-1-1} \textbf{Matrix identities.} If $A^2=I$, then $A^{-1}=A$. If $A^2=A$, then $A$ is idempotent and $A^n=A$ for every integer $n\ge1$. Also $A^{-1}=\operatorname{adj}(A)/\det A$ when $\det A\ne0$.
\item \label{formula-1-1-2} \textbf{Continuity, derivatives, and extrema.} A function is continuous at $a$ when $\lim_{x\to a}f(x)=f(a)$. A differentiable function can have an interior local extremum only where $f'(x)=0$. For $x>0$,
\begin{equation*}\tag{1.1.2.1}
\frac{d}{dx}x^x=x^x(1+\ln x).
\end{equation*}
\item \label{formula-1-1-3} \textbf{Trigonometric identity and integral.}
\begin{equation*}\tag{1.1.3.1}
1+\cos2x=2\cos^2x,\qquad
\int\sec x\,dx=\ln|\sec x+\tan x|+C.
\end{equation*}
Thus, on intervals where $\cos x>0$, $\int\frac{dx}{\sqrt{1+\cos2x}}=\frac1{\sqrt2}\ln|\sec x+\tan x|+C$.
\item \label{formula-1-1-4} \textbf{Area of an ellipse.} For $\frac{x^2}{a^2}+\frac{y^2}{b^2}=1$, the semi-axes are $a,b$ and the enclosed area is $\pi ab$.
\item \label{formula-1-1-5} \textbf{Linear programming.} For a linear objective $Z=ax+by$, find its maximum or minimum by evaluating $Z$ at each vertex of the feasible region.
\item \label{formula-1-1-6} \textbf{Conditional probability and Bayes' theorem.}
\begin{equation*}\tag{1.1.6.1}
P(A\mid B)=\frac{P(A\cap B)}{P(B)},\qquad
P(A\cup B)=P(A)+P(B)-P(A\cap B).
\end{equation*}
If $\{A_i\}$ partitions the sample space, $P(A_i\mid B)=\frac{P(B\mid A_i)P(A_i)}{\sum_jP(B\mid A_j)P(A_j)}$.
\item \label{formula-1-1-7} \textbf{Vectors.}
\begin{equation*}\tag{1.1.7.1}
\vec a\cdot\vec b=|\vec a||\vec b|\cos\theta,\qquad
\operatorname{Area}(\triangle ABC)=\frac12|\overrightarrow{AB}\times\overrightarrow{AC}|.
\end{equation*}
For unit vectors, $|\vec a-\vec b|^2=2-2\vec a\cdot\vec b$.
\item \label{formula-1-1-8} \textbf{Angle between lines.} If their direction vectors are $\vec d_1,\vec d_2$, then $\cos\theta=\frac{|\vec d_1\cdot\vec d_2|}{|\vec d_1||\vec d_2|}$ for the acute angle.
\item \label{formula-1-1-9} \textbf{Distance between skew lines and from a point to a line.}
\begin{equation*}\tag{1.1.9.1}
d=\frac{|(\vec a_2-\vec a_1)\cdot(\vec d_1\times\vec d_2)|}{|\vec d_1\times\vec d_2|}.
\end{equation*}
For a point $P$ and a line $\vec r=\vec a+t\vec d$, the distance is
\begin{equation*}\tag{1.1.9.2}
d(P,\ell)=\frac{|(\vec p-\vec a)\times\vec d|}{|\vec d|}.
\end{equation*}
\item \label{formula-1-1-10} \textbf{First-order linear differential equations.} For $x'+P(y)x=Q(y)$, the integrating factor is $\mu(y)=e^{\int P(y)\,dy}$ and $(\mu x)'=\mu Q$.
\item \label{formula-1-1-11} \textbf{Relations and functions.} If finite sets $A,B$ have $m,n$ elements, respectively, the number of relations from $A$ to $B$ is $2^{mn}$ and the number of functions from $A$ to $B$ is $n^m$.
\item \label{formula-1-1-12} \textbf{Linear systems in matrix form.} A system of linear equations can be written $AX=B$; if $\det A\ne0$, then $X=A^{-1}B$.

\item \label{formula-1-1-13} \textbf{Area of a triangle from coordinates.} For vertices $(x_1,y_1)$, $(x_2,y_2)$ and $(x_3,y_3)$, $[ABC]=\frac12|x_1(y_2-y_3)+x_2(y_3-y_1)+x_3(y_1-y_2)|$.
\end{enumerate}
